La integración de interfaces de chat asistidas por modelos de lenguaje en aplicaciones web introduce desafíos críticos relacionados con la latencia computacional. En los sistemas de mensajería convencionales, se emplean peticiones HTTP bajo el ciclo tradicional de solicitud y respuesta, en este modelo el cliente debe esperar a que el servidor genere la totalidad del mensaje antes de recibir cualquier dato. Sin embargo, debido a la naturaleza autorregresiva de los modelos de lenguaje de gran escala, la inferencia completa de una respuesta extensa puede demorar varios segundos, lo que resultaría en una mala experiencia de usuario debido a los tiempos de espera.

La solución adoptada por los líderes de la industria, tales como OpenAI (ChatGPT) o Anthropic (Claude) \cite{openai_streaming_2024}, consiste en la transmisión progresiva de datos o \textit{streaming}. Este mecanismo permite visualizar la respuesta fragmento a fragmento (\textit{token a token}) conforme el modelo la genera, optimizando drásticamente la latencia percibida. Técnicamente, este flujo de datos se implementa principalmente a través de \textbf{WebSockets} o de \textbf{Server-Sent Events (SSE)}. Mientras que el primer protocolo es más complejo al establecer una comunicación bidireccional y \textit{full-duplex} que requiere un cambio de protocolo (\textit{protocol upgrade}) y una gestión de estado más costosa en el servidor, el segundo destaca por su ligereza y su integración nativa con el estándar HTTP \cite{mackovivccomparative}.

En consecuencia, \textit{Server-Sent Events} se ha consolidado como el estandar en la arquitectura de servicios de IA generativa, ya que permite a las grandes plataformas mantener una infraestructura escalable y eficiente para flujos de datos unidireccionales, garantizando la fluidez de la interfaz sin la sobrecarga técnica de los canales bidireccionales.
